% ══════════════════════════════════════════════════════════════════════════════
\chapter{Lemmatisation et index inversé}
\label{chap:td3}
% ══════════════════════════════════════════════════════════════════════════════

La troisième étape apporte deux améliorations majeures à la représentation du corpus.
D'une part, la \textbf{lemmatisation} réduit chaque forme fléchie à sa forme canonique,
améliorant le rappel en regroupant sous un même terme des variantes morphologiques
(\texttt{voitures}, \texttt{voiture} → \texttt{voiture}). D'autre part, la construction
de l'\textbf{index inversé} — structure canonique de la RI selon le cours (chapitres 1
et 7) — fournit pour chaque terme la liste des documents qui le contiennent, permettant
de répondre en O(1) à toute requête booléenne.

\section{Comparaison des méthodes de lemmatisation}

Deux approches sont implémentées derrière une interface abstraite commune
\texttt{BaseLemmatizer}, permettant leur substitution transparente :

\begin{itemize}
  \item \textbf{SpaCy} (\texttt{fr\_core\_news\_sm}) — lemmatisation linguistique à
        base de modèle pré-entraîné ; produit des lemmes valides du dictionnaire
        (\texttt{innovations} $\to$ \texttt{innovation}).
  \item \textbf{Snowball} (NLTK) — racinisation par règles morphologiques suffixales ;
        produit des racines tronquées non nécessairement valides
        (\texttt{innovations} $\to$ \texttt{innov}).
\end{itemize}

Le critère de sélection automatique est le \textbf{taux de compression} $\tau_c =
|\mathcal{L}|/|\mathcal{V}|$, rapport entre le nombre de lemmes distincts et le nombre
de formes de surface distinctes. Un taux bas signifie un meilleur regroupement et donc
un meilleur rappel en RI.

\begin{table}[ht]
\centering
\caption{Comparaison quantitative des deux méthodes de lemmatisation sur 9\,843 formes}
\label{tab:lemma_compare}
\begin{tabular}{lrrrr}
\toprule
\textbf{Méthode} & \textbf{Formes} & \textbf{Lemmes/racines} & $\tau_c$ & \textbf{Regroupement} \\
\midrule
SpaCy    & 9\,843 & 7\,421 & 0,754 & 1,33$\times$ \\
Snowball & 9\,843 & 6\,192 & 0,629 & 1,59$\times$ \\
\bottomrule
\end{tabular}
\end{table}

Le tableau~\ref{tab:lemma_compare} montre que Snowball est sélectionné automatiquement
dans la plupart des exécutions grâce à son facteur de regroupement supérieur (1,59$\times$
contre 1,33$\times$). Sa principale limite est la lisibilité humaine des racines
produites (\texttt{technolog}, \texttt{recherch}), qui ne constituent pas des mots
valides du dictionnaire.

\section{Construction de l'index inversé}

La classe \texttt{InvertedIndexBuilder} parse le corpus XML une seule fois à
l'initialisation et expose quatre méthodes spécialisées selon le type de champ :

\begin{lstlisting}[caption={Index pondéré multi-champs : titre$\times$2 + texte$\times$1},
  label={lst:combined_index}]
def build_combined_index(
    self, fields: list[str], output_path: Path,
    weights: dict[str, float] | None = None,
) -> int:
    _weights = weights or {f: 1.0 for f in fields}
    index: dict[str, dict[str, float]] = defaultdict(lambda: defaultdict(float))
    for article_id, doc in self._iter_documents():
        for field in fields:
            el = doc.find(field)
            if el is None or not el.text:
                continue
            w = _weights.get(field, 1.0)
            for token in _tokenize(el.text):
                index[token][article_id] += w
    self._write_float_index(index, output_path)
    return len(index)
\end{lstlisting}

Le titre reçoit un poids double (\texttt{weights=\{"titre":~2.0,~"texte":~1.0\}}) car
il constitue un résumé thématique rédigé par l'auteur. Ce paramétrage s'inspire du
modèle BM25 et améliore le classement des documents dont le terme apparaît dans le
titre.

\section{Résultats}

Sept fichiers d'index sont produits, couvrant l'ensemble des axes de recherche utiles :
texte libre (titre, texte, combiné), facettes de navigation (rubrique, auteur, bulletin)
et filtre temporel (date par mois/année). La lemmatisation Snowball réduit le vocabulaire
de 56\,\% (14\,177 formes → 6\,192 racines), et un second filtrage TF-IDF sur les lemmes
élimine 47 lemmes grammaticaux supplémentaires (\texttt{être}, \texttt{avoir},
\texttt{pouvoir}\ldots) non capturés en TD2. La suite de 224 tests atteint 95,93\,\% de
couverture.
